\section{Boundary data, exterior sources, and contact coordinates}
\label{bdy:section}\label{sec:boundary}

We first replace the measurements by a family of interior source
solutions.  The boundary symbol determines jets of both the metric and
the connection, after a boundary-identity change of gauge.  These jets
allow a common exterior cap, on which the Green matrices agree.  Source
evaluations in that cap then identify points and fibers simultaneously.
The final part of the section prepares coordinates and harmonic frames
at a possible frontier of this identification.  The cap and
source-evaluation construction follows the scalar construction in
\cite[Section~2]{ScalarCompanion2026}; we give the changes needed for
unitary connections and matrix-valued evaluations.

\subsection{Dirichlet solutions and boundary jets}

We use positive metric densities throughout.  Write
\(L_{g,A}=d_A^{*g}d_A\), with positive principal symbol, and let
\(L_{g,A,D}\) denote its zero-Dirichlet realization.  Its quadratic form
is coercive on \(H^1_0(M;\mathbb C^2)\).  Indeed, the usual elliptic form
estimate and compactness reduce coercivity to the absence of a kernel.
A kernel element satisfies \(d_Au=0\); its norm is constant, and its
zero boundary trace makes it zero.  Smooth sources and boundary data
therefore have unique solutions smooth up to the boundary.  The same
facts hold on every smooth extension used below and on sufficiently
small coordinate balls; we use the standard smooth elliptic Dirichlet
theory of \cite[Chapter~5]{TaylorPDEI2011}.

We also use open-set unique continuation for systems with a smooth,
real scalar elliptic principal part and smooth matrix lower-order
terms; see \cite[Remark~3]{Aronszajn1957}.  The usual scalar local
Carleman estimate, summed over components, proves this statement: after conjugation by the scalar
weight, a matrix first-order term is bounded by a constant times
\(\|\nabla v\|+\tau\|v\|\), and is absorbed for large Carleman
parameter \(\tau\).  See \cite[Chapter~XXVIII]{HormanderIV2009} for the
underlying scalar estimates.  Boundary Cauchy uniqueness follows from
the same interior statement.  Extend the coefficients smoothly across
a boundary patch and extend a solution with zero trace and zero
covariant normal derivative by zero.  Green's formula leaves no
boundary distribution, so the extension solves the equation
weakly.  Interior regularity and unique continuation then give
vanishing near the patch and throughout the connected interior.

A unitary gauge \(V:M\to U(2)\) changes the connection to
\begin{equation}\label{bdy:gauge}
 A^V=V^{-1}AV+V^{-1}dV,\qquad
 d_{A^V}(V^{-1}u)=V^{-1}d_Au.
\end{equation}
A gauge equal to \(I\) on the boundary preserves the measured form.
In inward boundary-normal coordinates \((z,s)\), solve
\(\partial_sV=-A_sV\), \(V(z,0)=I\).  This gives \(A^V_s=0\).
The gauge can be made global: in a collar replace the time argument
of this solution by a smooth function equal to \(s\) in a smaller
collar and equal to zero near the inner edge, and set \(V=I\)
farther inside.  The ODE preserves unitarity.  Apply this construction
separately to \(A_1,A_2\), using gauges \(V_1,V_2\), and retain their
names for the final return to the original trivializations.  Until
then, \(A_j\) denotes the gauged connection.

\begin{lemma}[Boundary jets]\label{bdy:jets}
Suppose the local energy forms of \((g_1,A_1)\) and \((g_2,A_2)\)
agree on \(\Gamma\).  In their respective boundary-normal coordinates,
based on the same boundary coordinates, and in the normal gauges
just constructed, the full Taylor jets of \(g_1,A_1\) and
\(g_2,A_2\) agree at every point of \(\Gamma\).
\end{lemma}

\begin{proof}
Write \(g=ds^2+k(s,z)\), with \(k\) a tangential positive matrix,
and put \(\rho=(k^{ab}\xi_a\xi_b)^{1/2}\).  Green's formula identifies
the measured operator with the density
\(\nabla^A_\nu u\,dS_g\).  Relative to the coordinate density, its
principal symbol is
\[
             (\det k)^{1/2}\rho I.
\]
Its square determines \(C=(\det k)k^{-1}\).  Since the tangential
dimension is \(m=n-1\),
\[
 \det C=(\det k)^{m-1},\qquad
 k=(\det C)^{1/(m-1)}C^{-1}.
\]
Thus the boundary metrics and densities agree.  Division by their
common density gives equality of the ordinary covariant unit-normal
DN operators locally on \(\Gamma\).  Localization on both sides by
cutoffs supported in \(\Gamma\) gives equality of their full symbols.

We spell out the symbol calculation, extending the metric
factorization of \cite{LeeUhlmann1989} to the present scalar-principal
system.  Recovery of a connection together with the metric by this
even--odd symbol separation is also developed in
\cite[Theorem~1.1 and Section~3.3]{Gabdurakhmanov2021DN}; the
calculation below uses complex Hermitian fibers.  In normal gauge,
\[
 L_{g,A}=-\partial_s^2-\alpha\partial_s+P,
 \qquad \alpha=\tfrac12\operatorname{tr}(k^{-1}\partial_s k),
\]
where \(P\) is the connection Laplacian on each tangential slice.
With \(D_z=-i\partial_z\), its degree-one symbol consists of the
scalar metric term and \(-2ik^{ab}A_a\xi_b\).  Its degree-zero
symbol uses only slice coefficients and their tangential derivatives.
The local Dirichlet factorization has the form
\[
 L_{g,A}=(-\partial_s+B-\alpha)(\partial_s+B)
                 \pmod{\Psi^{-\infty}},
 \qquad
 b\mathbin{\#}b-\alpha b-\partial_s b\sim\sigma(P),
\]
where \(b\sim\rho I+b_0+b_{-1}+\cdots\) is a left symbol and
\[
 b\mathbin{\#}c\sim
 \sum_{\mu}\frac{(1/i)^{|\mu|}}{\mu!}
       (\partial_\xi^\mu b)(\partial_z^\mu c).
\]
The positive root \(\rho I\) selects the decaying Dirichlet branch.
The Poisson parametrix consequently identifies \(B(0)\), modulo a
smoothing operator, with the outward normal derivative.  This
factorization is valid for smooth coefficients: its principal root
is scalar and each subsequent matrix coefficient is solved by
scalar division by \(2\rho\).  The true inverse Dirichlet problem
changes only the local smoothing remainder.

The part of \(b_0\) involving the new jets is
\begin{equation}\label{bdy:first-jets}
 \frac14\left(\operatorname{tr}_k\partial_s k
       -\frac{(\partial_s k)(\xi^\sharp,\xi^\sharp)}{\rho^2}\right)I
       -\frac{ik^{ab}A_a\xi_b}{\rho},
 \qquad \xi^\sharp=k^{-1}\xi.
\end{equation}
This follows from
\(\partial_s\rho=-(\partial_s k)(\xi^\sharp,\xi^\sharp)/(2\rho)\)
and the degree-one factorization equation.  At degree \(1-j\),
\(j\ge2\), the new normal derivatives arise only in
\(\partial_s b_{2-j}/(2\rho)\).  Inductively their contribution is
\begin{equation}\label{bdy:highest-jets}
 \frac{1}{(2\rho)^{j-1}}
 \left[
 \frac14\left(\operatorname{tr}_k\partial_s^j k
       -\frac{(\partial_s^j k)(\xi^\sharp,\xi^\sharp)}{\rho^2}\right)I
       -\frac{ik^{ab}(\partial_s^{j-1}A_a)\xi_b}{\rho}
 \right].
\end{equation}
Every omitted term involves lower normal jets, with tangential
derivatives.  Differentiating \(\rho\), a raised index, or a
coefficient multiplying the preceding highest jet produces only
such lower-jet expressions.

Subtract the two symbol recursions after their lower jets have
been identified.  The metric term in
\eqref{bdy:highest-jets} is even in \(\xi\), whereas the connection
term is odd.  The odd term determines all components of the new
connection jet.  If the even term vanishes for a symmetric tensor
\(T\), then \(T(v,v)=\operatorname{tr}_k(T)\) for every unit vector
\(v\).  Polarization gives \(T=\operatorname{tr}_k(T)k\), and taking
trace gives \((m-1)\operatorname{tr}_k(T)=0\).  As \(m\ge2\),
\(T=0\).  Beginning with the recovered boundary metric, this proves
all normal jets by induction; their tangential derivatives recover
all mixed jets.  The remaining normal metric components and normal
connection component are fixed by the coordinate and gauge choices.
\end{proof}

\subsection{A common cap and its Green matrices}

Jet equality does not assert equality in an interior collar.  It
provides instead a common smooth extension on the exterior side of a
smaller patch.  The following variational argument turns that exterior
region into common interior data.

\begin{proposition}[Common exterior sources]\label{bdy:cap}
There are compact connected smooth extensions \(N_j\) of the original
manifold, with nonempty smooth boundary, and a common connected open
exterior cap \(E\subset N_j^\circ\), attached through a nonempty open
patch \(P\Subset\Gamma\), such that the extended metrics and
connections agree on \(E\).  For each
\(f\in C_c^\infty(E;\mathbb C^2)\), the solutions
\(L_{j,D}^{-1}f\) agree on \(E\).  The Dirichlet Green matrices,
normalized using metric volume, satisfy
\begin{equation}\label{bdy:common-green}
 G_1(x,y)=G_2(x,y)\quad(x,y\in E,\ x\ne y),
 \qquad G_j(x,y)=G_j(y,x)^*.
\end{equation}
\end{proposition}

\begin{proof}
Choose a boundary chart compactly contained in \(\Gamma\), and a
nonnegative smooth bump \(b\) whose support is compact in that chart
and whose positivity set \(P=\{b>0\}\) is nonempty and connected.
Using the separate normal collars, move the boundary from \(s=0\)
to the graph \(s=-b(z)\), with \(b\) small enough to remain inside
an extended collar.  The common open cap is
\[
                  E=\{(z,s):-b(z)<s<0\}.
\]
Extend the first coefficient fields smoothly to negative \(s\),
retaining positive definiteness of the metric and skew-Hermitian
values of the connection.  Use the same exterior extension for the
second pair.  Lemma~\ref{bdy:jets} makes the resulting pasting smooth
to all orders on a neighborhood of \(\operatorname{supp}b\).
Restricting it to \(s\ge-b(z)\) gives smoothness also at the flat
edges of the bump.  The bundles remain trivial in the extended
normal gauges.

For the energy comparison only, identify the extensions by the
identity on the original \(M\) and by cap coordinates on \(E\).
In collar charts this map and its inverse are continuous and
piecewise smooth with bounded first derivatives, and hence
bi-Lipschitz.  They identify the relevant \(H^1_0\) spaces.  Fix a
vector source \(f\) compactly supported in \(E\).  The solution
\(w_j=L_{j,D}^{-1}f\) is the unique minimizer of
\[
 \mathcal J_j(w)=\frac12\int_{N_j}|d_{A_j}w|_{g_j}^2\,dV_{g_j}
       -\operatorname{Re}\int_E\langle f,w\rangle\,dV_{g_j}
       \quad\text{on }H^1_0(N_j;\mathbb C^2).
\]
Its trace on the original boundary is smooth and supported in
\(\operatorname{supp}b\Subset\Gamma\), because the original and
extended boundaries coincide wherever \(b=0\).  Transport \(w_1\)
to \(N_2\), and inside the original \(M\) replace it by the
second harmonic extension of this trace.  The replacing difference
has zero trace on \(\partial M\), so its zero extension is in
\(H^1_0(N_2)\).  The competitor is therefore admissible even at the
narrowing edges of the cap.

On \(M\), the source is zero and \(w_1\) is the first harmonic
extension of its trace.  Equality of the measured quadratic forms
makes its interior energy equal to that of the replacement.  On
\(E\), both coefficients and the source pairing are unchanged.
Thus \(\min\mathcal J_2\le\min\mathcal J_1\).  Reversing the
argument gives equality, and uniqueness of the minimizer proves
\(w_1=w_2\) on \(E\).

Normalize the Green matrix by
\[
 L_{j,x}G_j(x,y)=\delta_y^{g_j}(x)I,\qquad
 G_j(\cdot,y)|_{\partial N_j}=0,\qquad
 L_{j,D}^{-1}f(x)=\int_{N_j}G_j(x,y)f(y)\,dV_{g_j}(y).
\]
The source densities on \(E\) coincide.  Varying the source gives
equality of the distribution kernels on \(E\times E\), hence
pointwise equality off the diagonal.  Self-adjointness of the
Dirichlet inverse gives the adjoint symmetry in
\eqref{bdy:common-green}.
\end{proof}

\subsection{Evaluations identify points and fibers}

Fix a nonempty open set \(U_0\Subset E\).  For
\(f\in C_c^\infty(U_0;\mathbb C^2)\), set
\begin{equation}\label{bdy:source-evaluations}
 w_{j,f}=L_{j,D}^{-1}f,\qquad I_j(p)f=w_{j,f}(p).
\end{equation}
Thus \(I_j(p)\) is a linear map from the common source space to the
fiber at \(p\).  The maps agree on \(E\).  This construction is
related to the source and Poisson embeddings of
\cite[Section~6.1 and Appendix~A]{LLS20}.  The following
point-distribution proof, in the spirit of
\cite[Chapter~III, Section~3]{Malgrange1956}, supplies the precise
system-valued separation needed here.

\begin{lemma}[Values and harmonic two-jets]\label{bdy:evaluations}
Let \(N\) be a compact connected smooth manifold of dimension
\(n\ge2\) with nonempty boundary, let \(L=d_A^{*g}d_A\), and let
\(U\Subset N^\circ\) be a nonempty open source set.
\begin{enumerate}
\item At each interior point, \(f\mapsto L_D^{-1}f(p)\) is onto
\(\mathbb C^2\).  At any two distinct interior points the two
values are jointly onto \(\mathbb C^2\oplus\mathbb C^2\).
\item If \(p\notin\overline U\), the second jets of \(L_D^{-1}f\)
at \(p\) fill exactly the linear space of jets satisfying
\(Lu(p)=0\).
\item Each evaluation is zero on \(\partial N\), and for fixed
\(f\) varies continuously up to that boundary.
\end{enumerate}
\end{lemma}

\begin{proof}
The last assertion is smooth Dirichlet regularity.  For the first
two, a complex-linear functional on values or jets is represented
by a dual-vector distribution \(T\), supported at one or two points.
Suppose it annihilates all source solutions, and define
\(v=(L_D^{-1})^tT\) by distributional transposition.  Then
\[
 \langle v,f\rangle=\langle T,L_D^{-1}f\rangle=0
                    \quad(f\in C_c^\infty(U;\mathbb C^2)),
 \qquad L^tv=T.
\]
Here the transpose and pairing are bilinear on the dual bundle;
no complex conjugation convention is needed.  Elliptic regularity
makes \(v\) smooth off the support of \(T\).  The interior with
one or two points removed is connected: a path can be diverted
around each point in a punctured coordinate ball when \(n\ge2\).
Unique continuation from the nonempty set obtained by removing
these points from \(U\) makes \(v\) supported at those points in
the interior.  If a point belongs to \(U\), the asserted initial
vanishing is distributional there as well.

At each support point a distribution is a finite sum of derivatives
of the point mass.  If its highest derivative order is \(k\), its
image under \(L^t\) has exact highest order \(k+2\).  Indeed, the
highest homogeneous coefficient vector polynomial is multiplied
by the scalar elliptic quadratic polynomial; this multiplication
is injective.  Consequently no nonzero order-zero distribution
can be the image of a point-supported distribution.  There are
therefore no value annihilators, proving both surjectivity claims.

For a two-jet annihilator at \(p\notin\overline U\), the image
\(T\) has order at most two.  Thus \(v=c\delta_p\), with \(c\)
a dual-fiber vector.  Its image is precisely the functional
\(u\mapsto c(Lu(p))\).  Conversely all such functionals
annihilate the source solutions.  Finite-dimensional duality gives
the exact stated jet space.  In particular, values and gradients
may be prescribed freely, since the Hessian trace can be chosen
to satisfy the equation.
\end{proof}

We can now define the correspondence whose continuation will prove
the theorem.  A \emph{match} consists of a local isometry
\(\Phi:W\to N_1^\circ\), with \(W\subset N_2^\circ\) connected,
open, and containing \(E\), and a smooth unitary bundle map \(J\)
from the second fiber at \(p\) to the first fiber at \(\Phi(p)\).
We require
\begin{equation}\label{bdy:matching}
 \Phi|_E=\operatorname{Id},\qquad J|_E=I,\qquad
 d_{\Phi^*A_1}(Ju)=Jd_{A_2}u,\qquad
 J(p)I_2(p)=I_1(\Phi(p)).
\end{equation}
In the connection identity the first connection is pulled back to
\(W\).  The identity on \(E\) is a match.

\begin{lemma}[The maximal match]\label{bdy:matches}
All matches agree on overlaps and are injective on base points.
Their union is a unique maximal match \((\Phi,J)\) on a connected
open set \(W\subset N_2^\circ\).  On this set,
\begin{equation}\label{bdy:matched-green}
 G_1(\Phi(p),\Phi(q))
       =J(p)G_2(p,q)J(q)^*\qquad(p,q\in W,\ p\ne q).
\end{equation}
\end{lemma}

\begin{proof}
If two possible images of the same point were distinct, their
evaluation identities would impose a fixed invertible linear
relation between the evaluations at two distinct points of
\(N_1^\circ\).  This contradicts joint surjectivity in
Lemma~\ref{bdy:evaluations}.  Once the image agrees, surjectivity
at the source point determines \(J\) uniquely.  The same argument
using the evaluations in \(N_2\) rules out two base points with
the same image.  Thus the maps glue on every overlap.  Their
union is connected because every domain contains \(E\), and
retains all the properties in \eqref{bdy:matching}.

For fixed \(p\), the evaluation identity and the common source
density give \eqref{bdy:matched-green} when \(q\in U_0\), away
from the pole.  As a function of \(q\), the adjoint of each side
solves the same pulled-back connection equation after the bundle
identification.  Injectivity of \(\Phi\) ensures that \(q=p\)
is the only possible pole.  Unique continuation on the connected
set \(W\setminus\{p\}\) proves the identity everywhere claimed.
\end{proof}

\subsection{Coordinates at a possible interior frontier}

It remains to show that the maximal domain is the entire interior.
We next prepare the local data needed to enlarge it.  The use of
harmonic \emph{frames} is essential: constant coefficient columns
in these frames solve the equation, even when the connection has
no nonzero parallel sections.

\begin{proposition}[Contact data]\label{bdy:contact}
If the maximal domain \(W\) has an interior frontier, there are a
point \(p\in\partial W\cap N_2^\circ\), a point \(q\in N_1^\circ\),
and charts about \(p,q\), both with coordinate range a neighborhood
of \(0\in\mathbb R^n\), with the following properties.
\begin{enumerate}
\item Both charts avoid \(\overline{U_0}\).  The old match is
coordinate identity on its domain in these charts.  That domain
contains a region \(x_n<\kappa(x')\), where \(\kappa\) is smooth,
\(\kappa(0)=0\), and \(d\kappa(0)=0\).
\item In smooth local unitary frames, the metrics, connection
matrices, and Green matrices agree on the known region (for the
Green matrices, on its square off the diagonal), and
\(g_1(0)=g_2(0)=I_n\).
\item There are invertible matrix functions \(F_j\) with harmonic
columns, chosen as corresponding source solutions.  In these
harmonic frames there are corresponding harmonic matrix functions
\(X_j^l\), \(1\le l\le n\), satisfying
\begin{equation}\label{bdy:trace-coordinates}
 \tfrac12\operatorname{Re}\operatorname{tr}X_j^l(x)=x^l,
 \qquad X_j^l(0)=0,
 \qquad \partial_iX_j^l(0)=\delta_i^l I.
\end{equation}
\item There are finitely many further corresponding harmonic
column functions with zero value and gradient at \(0\), whose
Hessians at \(0\) span the space of symmetric
\(\mathbb C^2\)-valued tensors \(H\) satisfying
\(\sum_i H_{ii}=0\).  All these paired functions, including the
frames, agree on the known region and have identical jets at
\(0\).
\end{enumerate}
\end{proposition}

\begin{proof}
Choose a small coordinate ball near an interior frontier and a
point of \(W\) close enough to that frontier that its distance
to the complement is attained inside the chart.  The Euclidean
ball with that center and radius lies in \(W\) and touches its
complement at an interior point \(p\).  Its boundary gives the
required smooth one-sided contact hypersurface.

If \(p_k\in W\) tends to \(p\), compactness of \(N_1\) and
\(U(2)\), using the fixed global trivializations, gives a
subsequence for which \(\Phi(p_k)\to q\) and \(J(p_k)\to J_0\).
Continuity of every source solution gives
\begin{equation}\label{bdy:limit-match}
                         J_0 I_2(p)=I_1(q).
\end{equation}
The point \(q\) cannot lie on the boundary, where the right side
would be zero, since \(I_2(p)\) is onto.  Two distinct possible
limits \(q\) contradict joint surjectivity on \(N_1\); at a fixed
\(q\), surjectivity of \(I_2(p)\) determines \(J_0\).  The full
limits therefore exist.  Also \(q\notin E\): otherwise its
identity match on \(E\) and \eqref{bdy:limit-match} would relate
the evaluations at the distinct points \(q,p\) of \(N_2\).
Both \(p\) and \(q\) consequently avoid \(\overline{U_0}\).

Choose two source solutions whose columns give an invertible
matrix \(F_1(q)\).  Their corresponding second-side columns give
\(F_2(p)\), invertible by \eqref{bdy:limit-match}.  Shrink the
charts so both matrices are invertible.  In a harmonic frame,
\(v=F_j^{-1}w\) satisfies an equation with scalar elliptic
principal part and no zeroth-order term, since every column of
\(F_j\) is harmonic.  Lemma~\ref{bdy:evaluations} is invariant
under this smooth frame change.  It therefore supplies matrices
\(X_1^l\) with value zero and gradients equal to any chosen real
coordinate covectors times \(I\).  Choose those covectors
independently, and let \(X_2^l\) be the corresponding source
matrices expressed in \(F_2\).

Set \(x^l=\frac12\operatorname{Re}\operatorname{tr}X_1^l\) and
\(y^l=\frac12\operatorname{Re}\operatorname{tr}X_2^l\).  The first
functions form coordinates at \(q\).  On matched points, the
source identity implies \(JF_2=F_1\circ\Phi\) and
\(X_2^l=X_1^l\circ\Phi\).  Thus the scalar coordinates and the
metric Gram matrices of their differentials agree under the
isometry.  Passing to \(p,q\), their Gram matrices agree and are
positive definite.  Hence the \(y^l\) are coordinates at \(p\).
Fix a target chart neighborhood \(U_1\) on which the \(x^l\)
are coordinates.  The full convergence of partner points permits
a source chart neighborhood \(U_2\) so small that
\(\Phi(W\cap U_2)\subset U_1\).  Choose a small coordinate ball
\(\Omega\) about zero contained in both \(x(U_1)\) and
\(y(U_2)\), and restrict the two charts to their inverse images
of \(\Omega\).  For a matched point in the restricted source
chart, \(x(\Phi(p'))=y(p')\in\Omega\); its partner therefore
lies in the restricted target chart.  Thus the old map is
coordinate identity throughout its domain in these charts.  The
matrix gradient identity at
\(q\) now reads \(\partial_iX_1^l(0)=\delta_i^lI\); equality on
the one-sided region and smoothness give the same identity for
\(X_2^l\).

Because the harmonic-frame equation has no zeroth-order term,
a jet with zero value and gradient is constrained only by
\(g_1^{ij}(0)H_{ij}=0\).  The jet lemma supplies a finite basis
of these Hessians and corresponding source pairs.  They agree
on the known region, so their complete jets agree at the contact
point.  The same conclusion holds for every smooth coefficient
or paired function already identified there.

Finally polar-orthonormalize the original harmonic frames:
write \(F_j=H_jS_j\), with \(H_j\) unitary and \(S_j\) positive
Hermitian.  On the matched set, \(JF_2=F_1\) implies
\(S_2=S_1\) and \(JH_2=H_1\).  In the unitary frames \(H_j\),
the old bundle identification is therefore \(I\); connection
matrices and Green matrices agree there.  We continue to write
\(F_j\) for the harmonic-frame matrices in these unitary frames.
A common real linear coordinate change makes the metric at zero
Euclidean and the contact hyperplane horizontal.  Apply the same
linear change to the list of \(X_j^l\); this preserves all
identities in \eqref{bdy:trace-coordinates}.  Reverse the final
normal coordinate if necessary to put the known side below its
graph.  In these coordinates the Hessian constraint is the one
in the statement.
\end{proof}

\subsection{Dilation of the contact data}\label{bdy:dilation}

The local analysis uses the preceding charts on a small physical
scale and a fixed coordinate range.  For \(\lambda>0\), set
\begin{equation}\label{bdy:dilation-formulas}
 \begin{aligned}
 g_{j,\lambda}(x)&=g_j(\lambda x),&
 A_{j,\lambda}(x)&=\lambda A_j(\lambda x),\\
 G_{j,\lambda}(x,y)&=\lambda^{n-2}G_j(\lambda x,\lambda y),&
 F_{j,\lambda}(x)&=F_j(\lambda x).
 \end{aligned}
\end{equation}
Here the metric normalization is \(\lambda^{-2}\) times the
pullback metric.  Thus the pulled-back connection Laplacian is
multiplied by \(\lambda^2\), and the factor
\(\lambda^{n-2}\) in the kernel exactly compensates for metric
volume.  In particular the displayed kernel still has unit
matrix point source for \(g_{j,\lambda}\).  Harmonic functions
and the harmonic frames dilate by composition, without a
multiplicative factor.

On every fixed bounded coordinate range, these metrics converge
smoothly to \(I_n\), the connection matrices converge smoothly
to zero, and all fixed finite orders of their coefficients and
frames remain bounded.  Source-free harmonic-frame operators
have zero zeroth-order terms; their ordinary-coordinate first-order
coefficients tend to zero.  The known side becomes
\(x_n<\lambda^{-1}\kappa(\lambda x')\), which contains every
fixed bounded region lying strictly below \(x_n=0\) for all
sufficiently small \(\lambda\).

We also have uniform kernel bounds on fixed separated sets.
An interior parametrix of order \(-2\) for the original smooth
scalar-principal system gives
\(|G_j(x,y)|\le C|x-y|^{2-n}\) in a sufficiently small fixed
chart, with a smooth remainder; interior estimates give the
corresponding bounds for derivatives away from the diagonal.
After \eqref{bdy:dilation-formulas}, every fixed positive lower
bound for \(|x-y|\) therefore gives bounds of any prescribed
finite order, uniformly for small \(\lambda\).  The local
constructions below first specify their bounded ranges and
positive separation margins and then restrict \(\lambda\).
They never demand a bound uniform as the separation tends to
zero from this initialization alone.

For the local argument, our convention for a normalized equation in
a harmonic frame is
\begin{equation}\label{bdy:harmonic-normalization}
 \mathcal L_jv=-c_jF_j^{-1}L_{g_j,A_j}(F_jv),\qquad c_j>0.
\end{equation}
The scalar multiplier acts outside the differential operator.
Its second-derivative coefficient matrix is \(c_jg_j^{-1}\),
and its zeroth-order term is zero.  This convention applies to
both the original and the dilated data; the functions \(c_j\)
will be chosen in Section~\ref{sec:spheres}.

The remaining local task is to continue the match through zero.
Sections~\ref{sec:product}--\ref{sec:matrix} will prove that the
metrics are conformal and that suitable normalizations
\eqref{bdy:harmonic-normalization} give equal operators on a
neighborhood of zero.  Section~\ref{glob:section} will then use
the vanishing zeroth-order terms to remove the conformal factor
and prove that \(F_1F_2^{-1}\) is a unitary connection
identification.  Unique continuation extends every source
evaluation, so the local identification enlarges the maximal match.
